\chapter{Implémentation et validation}
\label{chap:implementation}

\section{Environnement technique}

\begin{itemize}
\item Langages : Python 3.10, TypeScript/React.
\item Frameworks : PyTorch, FastAPI, Streamlit (prototypage).
\item Librairies : MNE, Scikit-learn, SHAP, Captum.
\item Conteneurisation : Docker, Kubernetes.
\end{itemize}

\section{Prétraitement des données}

Nous avons utilisé la base \textbf{Sleep-EDF} (version 2018) étendue, comprenant 196 sujets. Les étapes :
\begin{enumerate}
\item Chargement des fichiers EDF avec MNE.
\item Sélection des canaux :
  \begin{itemize}
    \item \textbf{2 canaux} : EEG Fpz-Cz et EEG Pz-Oz.
    \item \textbf{5 canaux} : les deux EEG + EOG horizontal + EMG submental + ECG (ou un canal supplémentaire selon la configuration).
  \end{itemize}
\item Filtrage passe-bande (0,5–35 Hz) et rééchantillonnage à 100 Hz.
\item Segmentation en épisodes de 30 secondes (3000 points).
\item Normalisation par épisode (z-score).
\item Division entraînement / validation / test : 70\% / 15\% / 15\% (au niveau des sujets).
\end{enumerate}

\section{Modèles entraînés}

Nous avons évalué trois architectures (LSTM, CNN, Transformer) sur deux configurations de canaux (2 et 5) et deux nombres de classes (3 : Éveil, NREM, REM ; et 5 : W, N1, N2, N3, REM). Soit un total de 12 modèles. Chaque modèle a été entraîné pendant 30 époques avec early stopping sur le kappa de Cohen.

Les hyperparamètres sont détaillés dans le tableau~\ref{tab:hyperparams}.

\begin{table}[H]
\centering
\caption{Hyperparamètres des modèles}
\begin{tabular}{|l|c|}
\hline
Paramètre & Valeur \\
\hline
Batch size & 64 \\
Optimiseur & Adam \\
Learning rate (LSTM/CNN) & 1e-3 \\
Learning rate (Transformer) & 1e-4 \\
Fonction de perte & CrossEntropyLoss \\
Dropout & 0,3 (LSTM/CNN), 0,1 (Transformer) \\
Nombre de couches LSTM & 2 (bidirectionnel) \\
Filtres CNN & 32, 64, 128 \\
Têtes d’attention (Transformer) & 4 \\
\hline
\end{tabular}
\end{table}

\section{Résultats expérimentaux}

\subsection{Performances globales}

Le tableau~\ref{tab:global_results} résume les meilleures précisions et kappa obtenus sur le jeu de test.

\begin{table}[H]
\centering
\caption{Résultats globaux des 12 modèles}
\begin{tabular}{|l|c|c|c|c|}
\hline
Modèle & Canaux & Classes & Accuracy (\%) & Kappa \\
\hline
LSTM & 5 & 3 & 90,86 & 0,8430 \\
CNN & 5 & 3 & \textbf{91,66} & \textbf{0,8527} \\
Transformer & 5 & 3 & 90,14 & 0,8260 \\
\hline
LSTM & 2 & 5 & 76,01 & 0,6793 \\
CNN & 2 & 5 & 79,90 & 0,7232 \\
Transformer & 2 & 5 & 77,15 & 0,6864 \\
\hline
LSTM & 2 & 3 & 87,58 & 0,7901 \\
CNN & 2 & 3 & 89,10 & 0,8095 \\
Transformer & 2 & 3 & 86,92 & 0,7714 \\
\hline
LSTM & 5 & 5 & 78,60 & 0,7131 \\
CNN & 5 & 5 & 82,78 & 0,7605 \\
Transformer & 5 & 5 & 81,55 & 0,7422 \\
\hline
\end{tabular}
\end{table}

Le meilleur modèle est le \textbf{CNN avec 5 canaux et 3 classes} (Éveil, NREM, REM) : précision de 91,66\% et kappa de 0,8527.

\subsection{Analyse détaillée du meilleur modèle}

Les matrices de confusion et rapports de classification sont présentés ci-dessous.

\begin{figure}[H]
\centering
\includegraphics[width=0.9\textwidth]{media/image3.png}  % à remplacer par votre fichier
\caption{Courbe d’apprentissage du CNN 5 canaux 3 classes}
\label{fig:cnn_5ch_3cl_curve}
\end{figure}

\begin{figure}[H]
\centering
\includegraphics[width=0.6\textwidth]{report/images/image4.png}
\caption{Matrice de confusion normalisée (CNN 5ch, 3 classes)}
\label{fig:cnn_5ch_3cl_cm}
\end{figure}

\begin{table}[H]
\centering
\caption{Rapport de classification détaillé (CNN 5ch, 3 classes)}
\begin{tabular}{|l|c|c|c|c|}
\hline
Classe & Précision & Rappel & F1-score & Support \\
\hline
Wake & 0,94 & 0,92 & 0,93 & 5990 \\
NREM & 0,92 & 0,94 & 0,93 & 11001 \\
REM & 0,84 & 0,79 & 0,82 & 2522 \\
\hline
Moyenne macro & 0,90 & 0,88 & 0,89 & 19513 \\
Moyenne pondérée & 0,92 & 0,92 & 0,92 & 19513 \\
\hline
\end{tabular}
\end{table}

On observe une très bonne discrimination entre Éveil et NREM, et une performance satisfaisante pour la classe REM (F1 = 0,82). Les erreurs principales se situent entre REM et NREM (confusion de 21\% des REM classés en NREM).

\subsection{Comparaison des architectures}

Le tableau~\ref{tab:comparison} montre que le CNN surpasse légèrement le LSTM et le Transformer pour la tâche à 3 classes. Pour la tâche à 5 classes, les performances chutent (~82\% max), ce qui justifie le choix de la granularité à 3 classes pour l’application clinique ciblée (détection des apnées).

\begin{table}[H]
\centering
\caption{Comparaison des architectures (5 canaux, 3 classes)}
\begin{tabular}{|l|c|c|}
\hline
Modèle & Accuracy (\%) & Kappa \\
\hline
LSTM & 90,86 & 0,8430 \\
CNN & \textbf{91,66} & \textbf{0,8527} \\
Transformer & 90,14 & 0,8260 \\
\hline
\end{tabular}
\end{table}

\section{Explicabilité : exemples avec SHAP}

Nous avons appliqué SHAP sur le meilleur modèle CNN pour expliquer les prédictions individuelles. La figure~\ref{fig:shap_rem} montre les valeurs SHAP pour un épisode correctement classé REM.

\begin{figure}[H]
\centering
\includegraphics[width=0.7\textwidth]{media/shap_rem_example.png} % à créer
\caption{Attributions SHAP pour un épisode REM}
\label{fig:shap_rem}
\end{figure}

Les échantillons autour de 250–500 ms (correspondant aux mouvements oculaires rapides) contribuent positivement à la classe REM, tandis que les zones de fort ralentissement (ondes lentes) contribuent négativement.

\section{Validation croisée et robustesse}

Une validation croisée à 5 plis sur les 196 sujets a confirmé la stabilité des performances : écart-type de l’accuracy inférieur à 0,8\% pour le modèle CNN 5 canaux 3 classes.

\section{Déploiement et performances temporelles}

L’inférence d’une nuit complète (environ 1000 épisodes) prend 8 minutes sur un GPU NVIDIA T4, soit 0,48 seconde par épisode. L’API FastAPI répond en moins de 2 secondes par requête. Le système conteneurisé (Docker) a été déployé sur un cluster Kubernetes de test.

\section{Synthèse des résultats}

\begin{itemize}
\item Le modèle \textbf{CNN à 5 canaux (EEG+EOG+EMG+ECG) en 3 classes (Éveil, NREM, REM)} atteint 91,66\% de précision, dépassant l’objectif des 90\% fixé dans le cahier des charges.
\item La version 5 classes reste utile pour la recherche, mais moins performante (82,78\% max).
\item Les explications SHAP sont cohérentes avec la physiologie du sommeil.
\item Le système est prêt pour une intégration clinique après validation complémentaire sur d’autres bases (MESA, SHHS).
\end{itemize}
